\section{«Bajar a la tierra»: la forma real del idealismo}

En la segunda mitad de la entrevista, el conductor insiste en la vision técnica. La visión que plantea Gao Jiyang (高继扬) es entrenar a los robots igual que se capacita a un empleado: con unas cuantas demostraciones y algunas prácticas por cuenta propia, que el robot complete tareas de forma estable y autónoma en su entorno. Para lograr esa experiencia, Xinghaitu (星海图) necesita tres componentes: el modelo base, las herramientas de postentrenamiento y el robot completo.

\begin{figure}[H]
\centering
\includegraphics[width=0.94\textwidth]{go-to-soil-robotics.png}
\caption{«Bajar a la tierra»: por qué emprender en robótica no es romántico.}
\end{figure}

\begin{knowledgebox}{Lectura de la figura: no ser romántico no significa carecer de idealismo}
A la izquierda, las aplicaciones de software y de modelos grandes tienden a concentrarse en el modelo, el crecimiento y la difusión; a la derecha, una empresa de robótica tiene que hacerse cargo del robot completo, la cadena de suministro, las instalaciones del cliente, el mantenimiento, la seguridad y el retorno de datos. La «tierra» del centro representa la fricción del mundo real. El idealismo de emprender en robótica no es un eslogan, sino mantener el objetivo después de enfrentar esas fricciones todos los días.
\end{knowledgebox}

\subsection{Cómo reconocer las promesas vacías}

Gao Jiyang considera que no es posible reconocer en un instante si una empresa de robótica hace promesas vacías; hay que observar un periodo: qué dijo el equipo hace uno o dos años, qué hizo durante el último año y si cumple dentro de un año. En la industria robótica, describir el futuro es necesario, porque los empleados, los inversionistas, los proveedores y los clientes necesitan creer en él; lo decisivo es si el equipo convierte de manera constante esa descripción del futuro en realidad.

\begin{importantbox}{La diferencia entre las promesas vacías y la narrativa estratégica}
Describir el futuro no es el problema; el problema es si ese futuro se desglosa en resultados por etapas, asignación de recursos, valor para el cliente y entregas verificables. Una narrativa que se cumple es estrategia; solo la que no se cumple es una promesa vacía.
\end{importantbox}

\subsection{El lobo: empuje y tenacidad}

Cuando le preguntan a qué animal se parece Xinghaitu, Gao Jiyang responde que lo más cercano es el lobo, aunque subraya que todas las empresas tienen algo de lobo. Aquí el «lobo» no alude a una cultura brusca, sino al empuje, la tenacidad y la determinación para alcanzar metas frente a las entregas a clientes y las fechas de lanzamiento. La intensidad de la competencia en la industria de la inteligencia corporizada exige que el equipo tenga idealismo y, al mismo tiempo, calcule con pragmatismo, todos los días, la relación entre inversión y resultados, y el valor a largo plazo.

\subsection{Resumen de la sección}

«Bajar a la tierra» es la metáfora industrial más importante de este episodio. Una empresa de robótica debe encontrar el equilibrio entre idealismo y pragmatismo: la visión tiene que llegar lejos, pero cada día hay que resolver problemas de hardware, cadena de suministro, clientes, datos y organización.
